% !TeX document-id = {bd0fbe59-c48a-4d54-be75-12fddd078f35}
% !TeX TXS-program:compile = txs:///xelatex
% !TeX TXS-program:bibliography = txs:///biber
% !TeX program = xelatex
% !TeX encoding = UTF-8

\documentclass[FM,BP]{tulthesis}
% ==============================================================================
% Šablona závěrečné práce Technické univerzity v Liberci (TUL)
% Verze šablony: 2.3 (tulthesis.cls)
%
% Autor šablony: Pavel Satrapa (http://www.nti.tul.cz/~satrapa/vyuka/latex-tul/)
% Autor komentářů, BibLaTeXu a ČSN ISO 690: Jan Koprnický (http://www.fm.tul.cz/personal/jan.koprnicky)
%
% NASTAVENÍ TŘÍDY DOKUMENTU:
% Fakulta:        FM  - Fakulta mechatroniky, informatiky a mezioborových studií
% Typ práce:      BP  - Bakalářská práce (pro bakalářskou práci je BP výchozí)
% Volitelné:      localfonts - aktivuje písma TUL z lokální složky ./pisma/ (pro XeLaTeX)
%                 fonts      - aktivuje systémově instalovaná písma TUL
%                 bw         - černobílý tisk (vhodné pro tištěnou verzi)
%                 sfheadings - bezserifové nadpisy kapitol
%                 twoside    - oboustranný tisk (pro finální vazbu)
% ==============================================================================

\newcommand{\verze}{2.3}

% Jazykové nastavení (XeLaTeX)
\usepackage{polyglossia}
\setdefaultlanguage{czech} % Primární jazyk: čeština
%\setdefaultlanguage{english} % Odkomentujte, pokud je celá práce v angličtině

% Matematické a typografické balíky
\usepackage{amsmath}
\usepackage{amssymb}
\usepackage{float}
\usepackage{makeidx}
\makeindex
\usepackage{dirtree}
\usepackage{tabularx}
\usepackage{tabularray}
\usepackage{booktabs}
\usepackage{graphicx}
\usepackage{xcolor}

\usepackage{xunicode}
\usepackage{xltxtra}

% Příkazy specifické pro šablonu
\newcommand{\argument}[1]{{\ttfamily\color{\tulcolor}#1}}
\newcommand{\argumentindex}[1]{\argument{#1}\index{#1}}
\newcommand{\prostredi}[1]{\argumentindex{#1}}
\newcommand{\prikazneindex}[1]{\argument{\textbackslash #1}}
\newcommand{\prikaz}[1]{\prikazneindex{#1}\index{#1@\textbackslash #1}}
\newenvironment{myquote}{\begin{list}{}{\setlength\leftmargin\parindent}\item[]}{\end{list}}
\newenvironment{listing}{\begin{myquote}\color{\tulcolor}}{\end{myquote}}
\sloppy

% ==============================================================================
% DEKLARACE METADAT PRO TITULNÍ STRÁNKU (IS STAG)
% ==============================================================================

% Název práce: {Česky}{Anglicky}
\TULtitle{Systém pro vzdálený monitoring FVE}%
{Remote Monitoring System for Solar Power Plants}

% Jméno autora
\TULauthor{Maksym Rudniev}

% Studijní program a obor (kódy a názvy dle zadání v IS STAG)
\TULprogramme{B0613A140005}{Informační technologie}{Information Technology}
\TULbranch{B0613A140005AI}{Aplikovaná informatika}{Applied Informatics}

% Vedoucí práce a případný konzultant
\TULsupervisor{Ing. Miloš Hernych}
%\TULconsultant{doc. RNDr. Jméno Příjmení, Ph.D.}

% Rok odevzdání
\TULyear{2026}

% ==============================================================================
% BIBLIOGRAFIE A CITACE (ČSN ISO 690 / BibLaTeX)
% ==============================================================================
\usepackage[
	backend=biber,
	style=iso-numeric,
	sorting=nyt,
	autolang=other,
	bibencoding=UTF8,
	maxcitenames=2,
	maxbibnames=3
]{biblatex}
\addbibresource{refTULTHESIS.bib}

% Hranaté závorky pro číslované citace v souladu s TUL
\DeclareFieldFormat{labelnumberwidth}{\mkbibbrackets{#1}}

% Oddělovače jmen autorů dle standardu
\DeclareDelimFormat{multinamedelim}{\addcomma\space}
\DeclareDelimFormat{finalnamedelim}{%
	\ifnumgreater{\value{liststop}}{2}{\finalandcomma}{}%
	\addspace\bibstring{and}\space}
\DeclareNameAlias{author}{family-given/given-family}

\usepackage{csquotes}
\urlstyle{same}

% Nastavení sazby zdrojového kódu (balík listings s podporou češtiny)
\usepackage{listings}
\makeatletter
\lst@InputCatcodes
\def\lst@DefEC{%
	\lst@CCECUse \lst@ProcessLetter
	^^80^^81^^82^^83^^84^^85^^86^^87^^88^^89^^8a^^8b^^8c^^8d^^8e^^8f%
	^^90^^91^^92^^93^^94^^95^^96^^97^^98^^99^^9a^^9b^^9c^^9d^^9e^^9f%
	^^a0^^a1^^a2^^a3^^a4^^a5^^a6^^a7^^a8^^a9^^aa^^ab^^ac^^ad^^ae^^af%
	^^b0^^b1^^b2^^b3^^b4^^b5^^b6^^b7^^b8^^b9^^ba^^bb^^bc^^bd^^be^^bf%
	^^c0^^c1^^c2^^c3^^c4^^c5^^c6^^c7^^c8^^c9^^ca^^cb^^cc^^cd^^ce^^cf%
	^^d0^^d1^^d2^^d3^^d4^^d5^^d6^^d7^^d8^^d9^^da^^db^^dc^^dd^^de^^df%
	^^e0^^e1^^e2^^e3^^e4^^e5^^e6^^e7^^e8^^e9^^ea^^eb^^ec^^ed^^ee^^ef%
	^^f0^^f1^^f2^^f3^^f4^^f5^^f6^^f7^^f8^^f9^^fa^^fb^^fc^^fd^^fe^^ff%
	% české znaky
	^^^^010c^^^^010d^^^^010e^^^^010f^^^^011a^^^^011b^^^^0147^^^^0148%
	^^^^0158^^^^0159^^^^0160^^^^0161^^^^0164^^^^0165%
	^^^^016e^^^^016f^^^^017d^^^^017e%
	^^00}
\lst@RestoreCatcodes
\makeatother

\lstset{
	basicstyle=\ttfamily\small,
	numbers=left,
	numberstyle=\tiny\color{gray},
	stepnumber=1,
	numbersep=8pt,
	backgroundcolor=\color{white},
	showspaces=false,
	showstringspaces=false,
	showtabs=false,
	frame=single,
	tabsize=4,
	captionpos=b,
	breaklines=true,
	breakatwhitespace=false,
	keywordstyle=\color{blue}\bfseries,
	commentstyle=\color{gray}\itshape,
	stringstyle=\color{teal}
}

% ==============================================================================
% HLAVNÍ DOKUMENT
% ==============================================================================
\begin{document}

	% Úvodní strany bakalářské práce:
	% \ThesisStart{male} vytvoří:
	% 1. Oficiální titulní stranu Bakalářské práce
	% 2. Místo pro vložení oficiálního zadání (vygenerovaného ze STAGu)
	% 3. Čestné prohlášení o samostatném vypracování bakalářské práce
	%
	% Po obdržení podepsaného zadání lze zadat přímo PDF soubor:
	% \ThesisStart{zadani-bp.pdf}
	\ThesisStart{male}

	% --------------------------------------------------------------------------
	% ABSTRAKT A KLÍČOVÁ SLOVA (ČESKY)
	% --------------------------------------------------------------------------
	\begin{abstractCZ}
	Tato bakalářská práce se zabývá návrhem, implementací a ověřením systému
	pro vzdálený monitoring fotovoltaických elektráren (FVE) a bateriových úložišť.
	Cílem práce je vytvoření spolehlivého řešení pro centralizovaný sběr, validaci,
	archivaci a vizualizaci telemetrických dat v reálném čase. V práci je navržena
	modulární backendová architektura využívající komunikační protokol MQTT pro příjem
	provozních dat ze střídače i odesílání řídicích stavů. Systém zahrnuje robustní
	validaci příchozích zpráv, dynamickou kompenzaci časového rozptylu (jitteru),
	perzistentní ukládání do optimalizované databáze SQLite v režimu WAL a integraci
	externích kontextových dat (spotové ceny elektřiny OTE a předpověď solárního osvitu).
	Pro zajištění vysoké provozní spolehlivosti je implementován mechanismus hlídacího
	psa (Watchdog) s automatickou detekcí výpadků komunikace a přechodem do bezpečnostního
	režimu (failsafe). Pro uživatele bylo vytvořeno přehledné webové rozhraní nabízející
	real-time vizualizaci okamžitých energetických a finančních toků, kapacity baterie,
	historických přehledů s filtrací a diagnostiky stavu spojení. Funkčnost celého systému
	byla experimentálně ověřena na reálných datech instalace Technické univerzity v Liberci.
	\end{abstractCZ}

	\begin{keywordsCZ}
	Fotovoltaická elektrárna, vzdálený monitoring, MQTT, databáze SQLite, telemetrie,
	vizualizace dat, Watchdog, internet věcí (IoT), spotový trh OTE.
	\end{keywordsCZ}

	\vspace{1.5cm}

	% --------------------------------------------------------------------------
	% ABSTRAKT A KLÍČOVÁ SLOVA (ANGLICKY)
	% --------------------------------------------------------------------------
	\begin{abstractEN}
	This bachelor thesis focuses on the design, implementation, and evaluation of a remote
	monitoring system for photovoltaic (PV) power plants and battery energy storage.
	The primary objective is to develop a reliable solution for centralized real-time telemetry
	acquisition, validation, archiving, and visualization. The thesis presents a modular backend
	architecture utilizing the MQTT protocol for ingesting operational inverter telemetry and
	dispatching control states. The system incorporates robust incoming data validation,
	dynamic packet jitter compensation, high-performance persistence in an optimized SQLite
	database operating in WAL mode, and integration of external operational context (day-ahead
	electricity spot prices and solar irradiance forecasts). To ensure high operational reliability,
	a communication watchdog mechanism is implemented with automated connection drop detection
	and fail-safe fallback. An interactive web-based interface provides real-time visualization
	of instantaneous energy and financial flows, battery state of charge, historical trends
	with filtering, and connection health diagnostics. The functionality of the complete system
	was experimentally verified using real-world operational data from the Technical University
	of Liberec installation.
	\end{abstractEN}

	\begin{keywordsEN}
	Photovoltaic power plant, remote monitoring, MQTT, SQLite database, telemetry,
	data visualization, Watchdog, Internet of Things (IoT), OTE spot market.
	\end{keywordsEN}

	\clearpage

	% --------------------------------------------------------------------------
	% PODĚKOVÁNÍ
	% --------------------------------------------------------------------------
	\begin{acknowledgement}
	Rád bych vyjádřil upřímné poděkování vedoucímu bakalářské práce,
	Ing. Miloši Hernychovi, za odborné vedení, cenné rady, vstřícný přístup
	a čas věnovaný konzultacím v průběhu řešení této práce. Dále děkuji
	své rodině a blízkým za stálou podporu během celého studia.
	\end{acknowledgement}

	% --------------------------------------------------------------------------
	% OBSAH A SEZNAMY
	% --------------------------------------------------------------------------
	\tableofcontents
	\clearpage

	%\listoffigures
	%\clearpage

	%\listoftables
	%\clearpage

	% --------------------------------------------------------------------------
	% SEZNAM ZKRATEK
	% --------------------------------------------------------------------------
	\begin{abbrList}
		\textbf{TUL}   & Technická univerzita v~Liberci \\
		\textbf{FM}    & Fakulta mechatroniky, informatiky a mezioborových studií TUL \\
		\textbf{FVE}   & Fotovoltaická elektrárna \\
		\textbf{BESS}  & Bateriové úložiště energie (Battery Energy Storage System) \\
		\textbf{SoC}   & Stav nabití akumulátoru (State of Charge) [\%] \\
		\textbf{OTE}   & Operátor trhu s~elektřinou, a.s. \\
		\textbf{LP}    & Lineární programování (Linear Programming) \\
		\textbf{MPC}   & Prediktivní řízení na bázi modelu (Model Predictive Control) \\
		\textbf{B2G}   & Dodávka energie z~baterie do sítě (Battery-to-Grid) \\
		\textbf{MQTT}  & Message Queuing Telemetry Transport \\
		\textbf{REST}  & Representational State Transfer (webové aplikační rozhraní) \\
		\textbf{GHI}   & Globální horizontální ozáření (Global Horizontal Irradiance) [$\text{W/m}^2$] \\
		\textbf{DNI}   & Přímé normálové ozáření (Direct Normal Irradiance) [$\text{W/m}^2$] \\
		\textbf{LFP}   & Lithium-železo-fosfátový akumulátor ($\text{LiFePO}_4$) \\
		\textbf{DPH}   & Daň z~přidané hodnoty (21\,\%) \\
		\textbf{PCC}   & Místo připojení k~distribuční síti (Point of Common Coupling) \\
	\end{abbrList}

	\clearpage

	% ==========================================================================
	% HLAVNÍ ČÁST PRÁCE (KAPITOLY)
	% ==========================================================================

	\chapter{Úvod}
	\label{chap:uvod}

	Využití domácích fotovoltaických elektráren (FVE) v kombinaci s bateriovým úložištěm energie (BESS) zažívá v České republice obrovský rozmach. Rostoucí počet domácností dnes opouští fixní tarify a přechází na dynamické spotové ceny elektřiny organizované operátorem trhu OTE. Na tomto trhu se cena silové elektřiny mění každou hodinu podle aktuální rovnováhy nabídky a poptávky, což otevírá zcela nové možnosti: v dobách drahé elektřiny omezit odběr ze sítě či energii výhodně prodávat, a naopak v hodinách přebytku nakupovat za minimální, nebo dokonce záporné ceny.

	Většina běžně instalovaných střídačů však dnes stále spoléhá na primitivní tovární logiku (režim \textit{Self-Consumption}). Ta funguje čistě reaktivně: jakmile ráno vyjde slunce, začne bez ohledu na budoucí ceny nabíjet baterii na maximum. Kolem poledne, kdy je solární výroba nejvyšší a spotové ceny často klesají do záporných hodnot, je však akumulátor již plně nabitý. Domácnost pak přebytečnou energii prodává se ztrátou, nebo musí výrobu nuceně omezovat (\textit{curtailment}). Naopak během drahých večerních špiček už v baterii často nezbývá dostatek energie.

	Právě tato neefektivita mě motivovala k tomu, abych se ve své práci zaměřil na návrh a ověření pokročilého prediktivního optimalizačního algoritmu pro řízení toků energie a bateriovou arbitráž (B2G -- \textit{Battery-to-Grid}).

	\section*{Čím jsem se v práci zabýval a co jsem zkoumal}
	Na počátku práce jsem se soustředil na detailní analýzu reálných provozních dat domácí FVE osazené hybridním střídačem a LiFePO$_4$ akumulátorem Fox-ESS o kapacitě 10,24\,kWh. Zkoumal jsem nejen chování spotřeby a solární výroby, ale především ekonomické a fyzikální zákonitosti, které v běžných teoretických modelech často chybí:
	\begin{itemize}
		\item \textbf{Reálnou strukturu českých tarifů:} Zisk z arbitráže nelze počítat pouze z čisté ceny OTE. Nákup ze sítě je zatížen regulovanými distribučními poplatky a 21\,\% DPH, zatímco při výkupu si obchodník strhává svou marži. Cenový spread tak musí pokrýt nemalé transakční náklady.
		\item \textbf{Ekonomiku opotřebení akumulátoru:} Každý cyklus způsobuje degradaci lithiových článků. Zkoumal jsem proto hranici rentability (\textit{Hurdle Rate}), aby zisk z arbitráže spolehlivě převýšil náklady na stárnutí baterie a ztráty vlivem účinnosti cyklu.
		\item \textbf{Český asymetrický metering:} V ČR je každá fáze měřena samostatně. Z analýzy telemetrie vyplynulo, že spotřeba domu je silně nevyvážená -- téměř polovinu energie odebírá fáze $L_3$ s elektrickým bojlerem.
		\item \textbf{Rizika spojená s předpovědí:} Zkoumal jsem, jak se vypořádat s chybami meteorologických předpovědí a jak zabránit nežádoucímu nočnímu přednabíjení ze sítě, které by přes den blokovalo kapacitu pro bezplatnou solární energii.
	\end{itemize}

	\section*{Návrh optimalizačního přístupu}
	Místo jednoduchých rozhodovacích pravidel jsem navrhl matematický model lineárního programování (LP/MILP) v rámci klouzavého prediktivního řízení (MPC) s 24hodinovým horizontem. Model v každém kroku zohledňuje předpověď osvitu, budoucí spotové ceny i stav akumulátoru a hledá globálně nejvýhodnější strategii -- kdy energii uskladnit, kdy pokrýt spotřebu domu a kdy ji komerčně exportovat do sítě. Detailnímu rozboru matematické formulace, ochranných podmínek a konkrétních strategií se věnuji v následujících kapitolách této práce.

	\section*{Čeho jsem dosáhl}
	Navržený algoritmus jsem implementoval v jazyce Python a podrobně otestoval na historických datech. Výsledky potvrdily, že prediktivní optimalizátor výrazně překonává standardní tovární řízení. Podařilo se eliminovat nákupy v drahých večerních špičkách, ochránit systém před ztrátami při záporných spotových cenách a zajistit bezpečný, ekonomicky rentabilní export energie. Dosažené výsledky prokázaly, že matematická optimalizace přináší pro moderní domácí mikrogridy zásadní finanční i provozní výhodu.


	\chapter{Optimalizační model a algoritmus bateriové arbitráže}
	\label{chap:optimalizator}

	V této kapitole je podrobně popsán matematický model a řídicí algoritmus
	vyvinutého energetického managementu (EMS -- \textit{Energy Management System})
	určeného pro komerční arbitráž na krátkodobém spotovém trhu s elektřinou.
	Algoritmus je navržen tak, aby autonomně rozhodoval o tocích energie v mikrogridu
	s fotovoltaickou elektrárnou a bateriovým úložištěm, maximalizoval finanční výnos
	a současně respektoval fyzikální limity technologie a specifika českého tržního prostředí.

	\section{Rámec prediktivního řízení a fyzikální model mikrosítě}
	\label{sec:mpc_ramec}

	Optimalizace je formulována v rámci klouzavého prediktivního řízení (MPC -- \textit{Model Predictive Control})
	s diskrétním časovým krokem $\Delta t = 1\,\text{h}$ a plánovacím horizontem $H = 24\,\text{hodin}$.
	Na počátku každé hodiny $k$ řešič vyhodnotí aktuální stav systému, aktualizuje předpovědi
	cen i počasí a vyřeší optimalizační úlohu pro celý horizont $t \in \{0, \dots, H-1\}$.
	Z vypočteného optimálního plánu se do reálného řízení aplikuje pouze akce pro první krok ($t = 0$),
	načež se v dalším časovém kroku celý proces s novými měřeními opakuje (princip uzavřené zpětné vazby).

	Uvažovaný mikrogrid se skládá ze čtyř základních entit:
	\begin{enumerate}
		\item \textbf{Fotovoltaická elektrárna (FVE):} Zdroj elektrické energie s predikovaným hodinovým profilem výroby $P_{\text{pv}}^{\text{pred}}(t) \ge 0$.
		\item \textbf{Nepružná zátěž objektu:} Domácí spotřeba $P_{\text{load}}^{\text{pred}}(t) \ge 0$, která je z pohledu optimalizátoru neovlivnitelná a musí být v každém okamžiku bezpodmínečně pokryta.
		\item \textbf{Bateriové úložiště (BESS):} Akumulátorový systém na bázi technologie $\text{LiFePO}_4$ (kalibrovaný podle modulární jednotky Fox-ESS HV2600 HS10.4 o jmenovité kapacitě $E_{\text{nom}} = 10{,}24\,\text{kWh}$) s maximálním trvalým výkonem nabíjení i vybíjení $P_{\text{bat}}^{\max} = 5{,}0\,\text{kW}$. Pro zajištění dlouhé životnosti ($\ge 6000$ cyklů) je provozní rozsah stavu nabití omezen na $\mathrm{SoC} \in [0{,}20; 1{,}00]$ jmenovité kapacity.
		\item \textbf{Bod připojení k distribuční síti (PCC):} Umožňuje obousměrnou výměnu výkonu -- nákup elektřiny $P_{\text{buy}}(t)$ a prodej $P_{\text{sell}}(t)$ až do limitu hlavního jističe $P_{\text{grid}}^{\max} = 10{,}0\,\text{kW}$.
	\end{enumerate}

	\section{Matematická formulace lineárního programování (MILP/LP)}
	\label{sec:lp_formulace}

	Pro každý časový krok $t \in \{0, \dots, H-1\}$ zavádíme vektor spojitých rozhodovacích proměnných:
	\begin{equation}
		x(t) = \big[ P_{\text{buy}}(t), \, P_{\text{sell}}(t), \, P_{\text{ch}}(t), \, P_{\text{dis}}(t), \, \mathrm{SoC}(t), \, P_{\text{curt}}(t) \big]^T \in \mathbb{R}^6_{\ge 0},
	\end{equation}
	kde $P_{\text{buy}}(t)$ je výkon odebíraný ze sítě, $P_{\text{sell}}(t)$ výkon dodávaný do sítě, $P_{\text{ch}}(t)$ nabíjecí výkon akumulátoru, $P_{\text{dis}}(t)$ vybíjecí výkon akumulátoru, $\mathrm{SoC}(t)$ úroveň energie v akumulátoru na konci kroku $t$ a $P_{\text{curt}}(t)$ řízené omezení solárního výkonu (\textit{curtailment}).

	\subsection{Účelová funkce (Objective Function)}
	Cílem optimalizace je minimalizace čistých finančních nákladů na provoz mikrogridu za horizont $H$, včetně nákladů na stárnutí akumulátoru a penalizace za zmařenou solární energii, s odečtením terminální hodnoty energie v baterii:
	\begin{equation}
		\label{eq:objective}
		\begin{split}
			\min \sum_{t=0}^{H-1} \Big( & P_{\text{buy}}(t) \cdot C_{\text{buy}}(t) - P_{\text{sell}}(t) \cdot C_{\text{sell}}(t) \\
			& + C_{\text{deg}} \cdot (P_{\text{ch}}(t) + P_{\text{dis}}(t)) + \epsilon \cdot P_{\text{curt}}(t) \Big) \Delta t - V_{\text{term}} \cdot \mathrm{SoC}(H-1).
		\end{split}
	\end{equation}
	Jednotlivé parametry účelové funkce představují:
	\begin{itemize}
		\item $C_{\text{buy}}(t)$ -- efektivní nákupní cena elektřiny ze sítě [Kč/kWh],
		\item $C_{\text{sell}}(t)$ -- efektivní výkupní cena elektřiny při exportu do sítě [Kč/kWh],
		\item $C_{\text{deg}}$ -- jednotkový degradační náklad na opotřebení bateriových článků ($C_{\text{deg}} = 0{,}20\,\text{Kč/kWh}$),
		\item $\epsilon$ -- malá numerická penalizace za nucené omezení přetoku ($\epsilon = 0{,}005\,\text{Kč/kWh}$), která zajišťuje, že curtailment nastane výhradně v situaci, kdy přebytečnou energii nelze využít ani prodat,
		\item $V_{\text{term}}$ -- terminální ocenění energie uložené v baterii na konci horizontu [Kč/kWh].
	\end{itemize}

	\subsection{Terminální hodnota energie \texorpdfstring{$V_{\text{term}}$}{V\_term}}
	V modelech s konečným horizontem bez terminální podmínky má řešič tendenci na konci intervalu ($t = H-1$) akumulátor kompletně vybít do sítě, protože budoucí hodnota energie po čase $H$ je v rovnici rovna nule. Aby k tomuto nežádoucímu jevu nedocházelo, stanovujeme terminální ocenění $V_{\text{term}}$ jako očekávanou úsporu z vytěsnění nákupu elektřiny v následujícím dni:
	\begin{equation}
		\label{eq:vterm}
		V_{\text{term}} = \max \left( 0{,}5; \, \min \Big( \bar{C}_{\text{buy}} \cdot \eta_{\text{dis}} - C_{\text{deg}}, \, 1{,}5 \cdot (\bar{C}_{\text{sell}} + C_{\text{deg}}) \Big) \right),
	\end{equation}
	kde $\bar{C}_{\text{buy}}$ a $\bar{C}_{\text{sell}}$ představují průměrné predikované ceny v horizontu. Tím je zajištěno, že baterie nevstoupí do dalšího dne uměle vyprázdněná.

	\subsection{Fyzikální a provozní omezení}
	Pro každý časový krok $t \in \{0, \dots, H-1\}$ musí být splněny následující omezující podmínky:

	\paragraph{1. Okamžitá uzlová bilance výkonu:}
	Zákon zachování energie v mikrogridu vyžaduje rovnost mezi zdroji a spotřebiči:
	\begin{equation}
		\label{eq:power_balance}
		P_{\text{buy}}(t) + P_{\text{pv}}^{\text{pred}}(t) + P_{\text{dis}}(t) = P_{\text{load}}^{\text{pred}}(t) + P_{\text{sell}}(t) + P_{\text{ch}}(t) + P_{\text{curt}}(t).
	\end{equation}

	\paragraph{2. Dynamika a kontinuita stavu nabití baterie:}
	Úroveň akumulované energie v čase $t$ závisí na předchozím stavu a tocích s uvážením účinností:
	\begin{equation}
		\label{eq:soc_continuity}
		\mathrm{SoC}(t) = \mathrm{SoC}(t-1) + \eta_{\text{ch}} \cdot P_{\text{ch}}(t) \cdot \Delta t - \frac{1}{\eta_{\text{dis}}} \cdot P_{\text{dis}}(t) \cdot \Delta t,
	\end{equation}
	přičemž pro $t = 0$ platí $\mathrm{SoC}(-1) = \mathrm{SoC}_{\text{curr}}$. Účinnosti nabíjení a vybíjení jsou nastaveny symetricky na $\eta_{\text{ch}} = \eta_{\text{dis}} = 0{,}975$, což odpovídá celkové cyklické účinnosti bateriového systému:
	\begin{equation}
		\eta_{\text{rt}} = \eta_{\text{ch}} \cdot \eta_{\text{dis}} = 0{,}975 \cdot 0{,}975 \approx 0{,}9506 \quad (95{,}06\,\%).
	\end{equation}

	\paragraph{3. Kapacitní a výkonové limity:}
	\begin{align}
		\mathrm{SoC}_{\min} &\le \mathrm{SoC}(t) \le \mathrm{SoC}_{\max}, \label{eq:soc_limits} \\
		0 &\le P_{\text{ch}}(t) \le P_{\text{bat}}^{\max}, \label{eq:pch_limits} \\
		0 &\le P_{\text{dis}}(t) \le P_{\text{bat}}^{\max}, \label{eq:pdis_limits} \\
		0 &\le P_{\text{buy}}(t) \le P_{\text{grid}}^{\max}, \label{eq:pbuy_limits} \\
		0 &\le P_{\text{sell}}(t) \le P_{\text{grid}}^{\max}, \label{eq:psell_limits} \\
		0 &\le P_{\text{curt}}(t) \le P_{\text{pv}}^{\text{pred}}(t), \label{eq:pcurt_limits}
	\end{align}
	kde $\mathrm{SoC}_{\min} = 0{,}20 \cdot 10{,}24 = 2{,}048\,\text{kWh}$, $\mathrm{SoC}_{\max} = 10{,}24\,\text{kWh}$, $P_{\text{bat}}^{\max} = 5{,}0\,\text{kW}$ a $P_{\text{grid}}^{\max} = 10{,}0\,\text{kW}$.

	\section{Tarifní a regulatorní model českého trhu}
	\label{sec:tarify}

	Správná kalkulace efektivních cen je pro ekonomickou úspěšnost arbitráže zcela zásadní. V České republice je nákup i prodej na spotovém trhu zatížen asymetrickými poplatky a daněmi:

	\paragraph{Efektivní nákupní cena $C_{\text{buy}}(t)$:}
	Zahrnuje hodinovou spotovou cenu OTE $P_{\text{spot}}(t)$, regulovanou složku za distribuci, podporu obnovitelných zdrojů (POZE), systémové služby ČEPS a daň z přidané hodnoty (DPH 21\,\%):
	\begin{equation}
		\label{eq:cbuy}
		C_{\text{buy}}(t) = \big( P_{\text{spot}}(t) + \text{Fee}_{\text{import}} \big) \cdot (1 + \text{DPH}),
	\end{equation}
	kde distribuční poplatek $\text{Fee}_{\text{import}} = 1{,}80\,\text{Kč/kWh}$ a $\text{DPH} = 0{,}21$.

	\paragraph{Efektivní prodejní cena $C_{\text{sell}}(t)$:}
	Při dodávce energie do sítě domácnost obdrží spotovou cenu sníženou o marži obchodníka (agregátora):
	\begin{equation}
		\label{eq:csell}
		C_{\text{sell}}(t) = P_{\text{spot}}(t) - \text{Fee}_{\text{export}},
	\end{equation}
	kde srážka obchodníka činí $\text{Fee}_{\text{export}} = 0{,}45\,\text{Kč/kWh}$.

	\section{Teoretická odvození a matematické důkazy}
	\label{sec:dukazy}

	V této části jsou exaktně odvozeny klíčové podmínky, které odlišují funkční komerční arbitráž od naivních přístupů.

	\subsection{Podmínka rentability arbitráže (Hurdle Rate Theorem)}

	\paragraph{Věta 1 (Hurdle Rate):}
	\textit{Nechť systém nakoupí ze sítě energii v čase $t_1$ a prodá ji zpět do sítě v čase $t_2$. Tato tržní operace je striktně zisková tehdy a jen tehdy, pokud efektivní prodejní cena v čase $t_2$ splňuje nerovnost:}
	\begin{equation}
		\label{eq:hurdle_rate}
		C_{\text{sell}}(t_2) > \frac{C_{\text{buy}}(t_1)}{\eta_{\text{rt}}} + C_{\text{deg}} \cdot \left( 1 + \frac{1}{\eta_{\text{rt}}} \right).
	\end{equation}

	\paragraph{Důkaz:}
	Uvažujme dodávku energie $E_{\text{net}} = 1\,\text{kWh}$ do distribuční sítě v čase $t_2$.
	K tomu, aby střídač dodal do sítě $1\,\text{kWh}$, musí z chemické kapacity akumulátoru odebrat energii:
	\begin{equation}
		E_{\text{bat,dis}} = \frac{E_{\text{net}}}{\eta_{\text{dis}}} = \frac{1}{\eta_{\text{dis}}}\,\text{kWh}.
	\end{equation}
	Tato energie musela být dříve do akumulátoru uložena v čase $t_1$. K uložení energie $E_{\text{bat,dis}}$ do baterie bylo nutné ze sítě nakoupit:
	\begin{equation}
		E_{\text{grid,in}} = \frac{E_{\text{bat,dis}}}{\eta_{\text{ch}}} = \frac{1}{\eta_{\text{ch}} \cdot \eta_{\text{dis}}} = \frac{1}{\eta_{\text{rt}}}\,\text{kWh}.
	\end{equation}
	Celkový finanční náklad na nákup elektřiny činí:
	\begin{equation}
		N_{\text{nakup}} = E_{\text{grid,in}} \cdot C_{\text{buy}}(t_1) = \frac{C_{\text{buy}}(t_1)}{\eta_{\text{rt}}}.
	\end{equation}
	Dále operace způsobuje degradaci článků při nabíjení i vybíjení. Celkový degradační náklad je:
	\begin{equation}
		N_{\text{deg}} = C_{\text{deg}} \cdot (E_{\text{grid,in}} + E_{\text{net}}) = C_{\text{deg}} \cdot \left( \frac{1}{\eta_{\text{rt}}} + 1 \right).
	\end{equation}
	Příjem z prodeje energie činí $T_{\text{prodej}} = E_{\text{net}} \cdot C_{\text{sell}}(t_2) = C_{\text{sell}}(t_2)$.
	Čistý zisk $\Pi = T_{\text{prodej}} - (N_{\text{nakup}} + N_{\text{deg}})$ je kladný právě tehdy, když:
	\begin{equation}
		C_{\text{sell}}(t_2) - \frac{C_{\text{buy}}(t_1)}{\eta_{\text{rt}}} - C_{\text{deg}} \cdot \left( 1 + \frac{1}{\eta_{\text{rt}}} \right) > 0,
	\end{equation}
	což přímo dokazuje vztah~\eqref{eq:hurdle_rate}. \hfill $\blacksquare$

	\paragraph{Kvantitativní důsledek pro český trh:}
	Dosadíme-li typické parametry: nákupní spotovou cenu $P_{\text{spot}}(t_1) = 1{,}00\,\text{Kč/kWh}$, dostaneme:
	\begin{equation}
		C_{\text{buy}}(t_1) = (1{,}00 + 1{,}80) \cdot 1{,}21 = 3{,}388\,\text{Kč/kWh}.
	\end{equation}
	Při $\eta_{\text{rt}} = 0{,}9506$ a $C_{\text{deg}} = 0{,}20\,\text{Kč/kWh}$ činí minimální efektivní prodejní cena:
	\begin{equation}
		C_{\text{sell}}(t_2) > \frac{3{,}388}{0{,}9506} + 0{,}20 \cdot \left( 1 + \frac{1}{0{,}9506} \right) = 3{,}564 + 0{,}410 = 3{,}974\,\text{Kč/kWh}.
	\end{equation}
	Jelikož $P_{\text{spot}}(t_2) = C_{\text{sell}}(t_2) + \text{Fee}_{\text{export}} = 3{,}974 + 0{,}45 = 4{,}424\,\text{Kč/kWh}$, nákup za $1{,}00\,\text{Kč}$ se vyplatí prodat až při spotové ceně přesahující $4{,}42\,\text{Kč/kWh}$!
	Tento výpočet exaktně dokazuje, proč naivní arbitráž bez Hurdle Rate v praxi generuje setrvalou finanční ztrátu.

	\subsection{Vytěsnění solární energie a pravidlo Zero Grid Pre-Charge}

	\paragraph{Věta 2 (Zero Grid Pre-Charge):}
	\textit{Nechť je předpovězen slunečný den, pro který očekávaný přebytek solární výroby nad spotřebou domu $S_{\text{pv}}^{\text{safe}}$ převyšuje volnou kapacitu baterie: $S_{\text{pv}}^{\text{safe}} \ge \mathrm{SoC}_{\max} - \mathrm{SoC}(t)$. Pak jakékoliv noční nabíjení baterie ze sítě ($P_{\text{ch}} > 0, P_{\text{buy}} > 0$) v hodinách před solárním vrcholem vede k čisté finanční ztrátě oproti nulovému nabíjení.}

	\paragraph{Důkaz:}
	Každá $1\,\text{kWh}$ energie uložená do baterie ze sítě v nočních hodinách za náklad $C_{\text{buy}}^{\text{night}} > 0$ obsadí $1\,\text{kWh}$ kapacity úložiště.
	V průběhu dne nastane situace, kdy solární panely vygenerují přebytek energie s nulovým marginálním výrobním nákladem ($C_{\text{pv}} = 0$). Protože je však kapacita baterie obsazena dříve nakoupenou elektřinou, musí být tato solární energie buď:
	\begin{enumerate}
		\item nuceně zmařena ($P_{\text{curt}} > 0$), což přináší finanční výnos $0\,\text{Kč}$, přičemž noční náklad $C_{\text{buy}}^{\text{night}}$ zůstal uhrazen, nebo
		\item prodána do sítě za polední cenu $C_{\text{sell}}^{\text{day}}$.
	\end{enumerate}
	Vzhledem k distribučním poplatkům a DPH však platí:
	\begin{equation}
		C_{\text{sell}}^{\text{day}} = P_{\text{spot}}^{\text{day}} - 0{,}45 < (P_{\text{spot}}^{\text{night}} + 1{,}80) \cdot 1{,}21 = C_{\text{buy}}^{\text{night}},
	\end{equation}
	což je v poledních hodinách slunečných dnů (kdy spotové ceny navíc klesají vlivem celostátní solární nadprodukce) splněno téměř univerzálně.
	Přednabití ze sítě tedy vytěsnilo bezplatnou solární energii a nahradilo ji drahou energií ze sítě. Optimální strategie proto pro noční hodiny před slunečným dnem striktně stanovuje:
	\begin{equation}
		P_{\text{ch}}(t) \le P_{\text{pv}}^{\text{pred}}(t) \quad \forall t \in \{0, \dots, 5\},
	\end{equation}
	čímž je noční odběr do baterie zakázán. \hfill $\blacksquare$

	\subsection{Záporné ceny a eliminace ztrát (Dynamic Curtailment)}

	\paragraph{Věta 3 (Ochrana proti záporným cenám):}
	\textit{Při výskytu záporné efektivní výkupní ceny $C_{\text{sell}}(t) < 0$ optimální řešení úlohy~\eqref{eq:objective} vynucuje $P_{\text{sell}}(t) = 0$ a přebytečnou solární energii přesměrovává do proměnné $P_{\text{curt}}(t)$.}

	\paragraph{Důkaz:}
	V účelové funkci~\eqref{eq:objective} vystupuje člen $-P_{\text{sell}}(t) \cdot C_{\text{sell}}(t)$.
	Je-li $C_{\text{sell}}(t) < 0$, platí $-C_{\text{sell}}(t) > 0$, a prodej do sítě tedy působí jako přímá pokuta zvyšující hodnotu minimalizované účelové funkce.
	Pokud je akumulátor plně nabit ($\mathrm{SoC} = \mathrm{SoC}_{\max}$), platí $P_{\text{ch}}(t) = 0$. V bilanční rovnici~\eqref{eq:power_balance} potom vystupuje přebytek:
	\begin{equation}
		P_{\text{pv}}^{\text{pred}}(t) - P_{\text{load}}^{\text{pred}}(t) = P_{\text{sell}}(t) + P_{\text{curt}}(t).
	\end{equation}
	Jelikož penalizace za curtailment je nastavena na $\epsilon = 0{,}005\,\text{Kč/kWh}$, zatímco pokuta za prodej při záporné ceně činí $|C_{\text{sell}}(t)| \ge 0{,}45\,\text{Kč/kWh}$, řešič lineárního programování lineárně preferuje $P_{\text{curt}}(t)$ před $P_{\text{sell}}(t)$. Tím je zajištěno $P_{\text{sell}}(t) = 0$ bez narušení bilanční rovnosti a bez pádu řešiče do nepřípustnosti (\textit{infeasibility}). \hfill $\blacksquare$

	\section{Strategie řízení a robustnost proti nejistotě}
	\label{sec:strategie}

	Pro praktické nasazení byl základní LP model rozšířen o robustní bezpečnostní vrstvy:

	\subsection{Konzervativní bezpečnostní diskonty předpovědí}
	Predikce počasí i spotřeby nesou inherentní nejistotu. Aby domácnost nebyla vystavena riziku nedostatku energie při náhlém zakalení oblohy, pracuje optimalizátor s konzervativně upravenými profily:
	\begin{align}
		P_{\text{pv}}^{\text{safe}}(t) &= 0{,}80 \cdot P_{\text{pv}}^{\text{pred}}(t) \quad (\text{20\% bezpečnostní diskont na solární výrobu}), \\
		P_{\text{load}}^{\text{safe}}(t) &= 1{,}15 \cdot P_{\text{load}}^{\text{pred}}(t) \quad (\text{15\% bezpečnostní přirážka na spotřebu domu}).
	\end{align}
	Výpočet očekávaného slunečního přebytku potom probíhá jako:
	\begin{equation}
		S_{\text{pv}}^{\text{safe}} = \sum_{t=0}^{H-1} \max \left( 0{,}0; \, P_{\text{pv}}^{\text{safe}}(t) - P_{\text{load}}^{\text{safe}}(t) \right).
	\end{equation}

	\subsection{Adaptivní držení energie pod cenové špičky (Adaptive Peak Preservation)}
	Během odpoledních hodin (13:00--18:00) dochází často k poklesu solární výroby. Naivní řízení v této době začne okamžitě vybíjet akumulátor na pokrytí drobných spotřebičů v domě. Algoritmus B2G monitoruje, zda v podvečerních hodinách (19:00--22:00) nenastává výrazná cenová špička ($P_{\text{spot}} \ge 6{,}00\,\text{Kč/kWh}$). Pokud ano, zablokuje předčasné vybíjení a akumulovanou energii rezervuje na večer, kdy buď pokryje vysoké odběry, nebo ji se značným ziskem vyexportuje do sítě.

	\subsection{Pravidlo Battery First v reálném čase}
	Plán vypočtený MPC představuje hodinový rámec. V subsekundovém reálném čase však v domě dochází k dynamickým spínáním spotřebičů (např. indukční deska, rychlovarná konvice). Pravidlo \textit{Battery First} funguje jako okamžitá zpětná vazba: jakýkoliv okamžitý výkonový deficit je primárně kompenzován z akumulátoru až do limitu $P_{\text{bat}}^{\max}$, čímž je zamezeno nechtěným odběrům ze sítě v drahých tarifech.

	\section{Experimentální benchmarky a vyhodnocení na reálných datech}
	\label{sec:benchmarky}

	Navržený optimalizační algoritmus byl podroben rozsáhlému experimentálnímu testování na reálných datech z provozu domácí FVE za období od května 2025 do září 2026 (celkem 17 měsíců, 11\,686 hodin nepřetržitého provozu). 

	Jako srovnávací báze (\textit{Baseline}) byla použita standardní tovární logika \textit{Self-Consumption} (naivní nabíjení při solárním přebytku, okamžité vybíjení při deficitu, neřízený export a žádná reakce na spotové ceny).

	\subsection{Souhrnné měsíční výsledky}
	Tabulka~\ref{tab:benchmark_results} shrnuje dosažené energetické a ekonomické ukazatele za celé testovací období.

	\begin{table}[htbp]
		\centering
		\small
		\caption{Měsíční srovnání referenční strategie (Naive) a navrženého algoritmu Arbitrage B2G (reálná data květen 2025 -- září 2026)}
		\label{tab:benchmark_results}
		\resizebox{\textwidth}{!}{%
		\begin{tabular}{lrrrrrrrrr}
			\toprule
			\textbf{Měsíc} & \textbf{Hodin} & \textbf{Výroba FVE} & \textbf{Spotřeba} & \textbf{Výsledek Naive} & \textbf{Výsledek B2G} & \textbf{Přínos B2G} & \textbf{B2G Export} & \textbf{Curtailment} & \textbf{Cykly} \\
			& [h] & [kWh] & [kWh] & [Kč] & [Kč] & [Kč] & [kWh] & [kWh] & [bat] \\
			\midrule
			2025-05 & 177 & 234{,}2 & 304{,}2 & -778{,}66 & -501{,}59 & \textbf{+277{,}06} & 24{,}1 & 28{,}8 & 8{,}1 \\
			2025-06 & 720 & 716{,}9 & 694{,}0 & -1\,440{,}88 & -739{,}43 & \textbf{+701{,}45} & 126{,}4 & 92{,}5 & 24{,}2 \\
			2025-07 & 744 & 322{,}6 & 240{,}7 & -217{,}52 & +90{,}75 & \textbf{+308{,}27} & 108{,}2 & 1{,}6 & 16{,}1 \\
			2025-08 & 744 & 377{,}6 & 348{,}5 & -360{,}95 & -33{,}59 & \textbf{+327{,}36} & 83{,}1 & 13{,}3 & 17{,}1 \\
			2025-09 & 720 & 223{,}9 & 219{,}1 & -285{,}72 & +37{,}81 & \textbf{+323{,}53} & 102{,}2 & 7{,}2 & 15{,}5 \\
			2025-10 & 744 & 126{,}5 & 194{,}1 & -612{,}19 & -368{,}98 & \textbf{+243{,}21} & 6{,}4 & 0{,}0 & 5{,}7 \\
			2025-11 & 720 & 91{,}3 & 177{,}4 & -678{,}61 & -476{,}55 & \textbf{+202{,}06} & 0{,}0 & 0{,}0 & 3{,}2 \\
			2025-12 & 744 & 50{,}3 & 149{,}6 & -743{,}75 & -522{,}70 & \textbf{+221{,}06} & 0{,}0 & 0{,}0 & 7{,}6 \\
			2026-01 & 744 & 43{,}0 & 135{,}2 & -730{,}90 & -516{,}67 & \textbf{+214{,}24} & 12{,}0 & 0{,}0 & 9{,}3 \\
			2026-02 & 672 & 95{,}2 & 163{,}8 & -567{,}63 & -389{,}75 & \textbf{+177{,}88} & 2{,}3 & 0{,}0 & 7{,}5 \\
			2026-03 & 744 & 250{,}1 & 311{,}2 & -584{,}31 & -319{,}83 & \textbf{+264{,}48} & 18{,}5 & 0{,}1 & 10{,}9 \\
			2026-04 & 720 & 370{,}3 & 403{,}0 & -474{,}97 & +62{,}30 & \textbf{+537{,}27} & 81{,}6 & 5{,}1 & 20{,}5 \\
			2026-05 & 744 & 382{,}4 & 395{,}2 & -527{,}48 & -105{,}76 & \textbf{+421{,}72} & 66{,}1 & 24{,}0 & 17{,}0 \\
			2026-06 & 720 & 344{,}8 & 237{,}2 & -103{,}39 & +461{,}65 & \textbf{+565{,}04} & 146{,}5 & 6{,}2 & 19{,}4 \\
			2026-07 & 744 & 349{,}2 & 186{,}1 & -121{,}34 & +349{,}43 & \textbf{+470{,}77} & 169{,}4 & 17{,}0 & 21{,}8 \\
			2026-08 & 744 & 299{,}2 & 184{,}2 & -57{,}01 & +348{,}35 & \textbf{+405{,}35} & 139{,}5 & 3{,}2 & 16{,}7 \\
			2026-09 & 541 & 186{,}3 & 169{,}1 & -231{,}51 & +124{,}40 & \textbf{+355{,}90} & 87{,}8 & 0{,}0 & 13{,}5 \\
			\midrule
			\textbf{Celkem} & \textbf{11\,686} & \textbf{4\,463{,}8} & \textbf{4\,512{,}6} & \textbf{-8\,516{,}82} & \textbf{-2\,500{,}16} & \textbf{+6\,016{,}65} & \textbf{1\,174{,}1} & \textbf{199{,}0} & \textbf{234{,}1} \\
			\bottomrule
		\end{tabular}%
		}
	\end{table}

	\subsection{Analýza a diskuse výsledků}
	Z naměřených a nasimulovaných dat vyplývají následující klíčové závěry:

	\paragraph{1. Celkový ekonomický přínos:}
	Zatímco referenční strategie vygenerovala za sledované období celkový náklad na energii ve výši $-8\,516{,}82\,\text{Kč}$, navržený algoritmus snížil celkový náklad na $-2\,500{,}16\,\text{Kč}$. Čistá finanční úspora činí \textbf{+6\,016{,}65\,Kč}, což představuje \textbf{snížení celkových výdajů na elektřinu o více než 70\,\%}.

	\paragraph{2. Sezónní chování a komerční export:}
	\begin{itemize}
		\item \textbf{Letní měsíce (červen až srpen):} V letním období se systém stává plně ziskovým. Například v červnu 2026 dosáhl čistého měsíčního zisku $+461{,}65\,\text{Kč}$ a v červenci $+349{,}43\,\text{Kč}$ (oproti záporné bilanci u Naive). Celkově bylo v lukrativních špičkách prodáno $1\,174{,}1\,\text{kWh}$ elektřiny.
		\item \textbf{Zimní měsíce (listopad až únor):} V zimě, kdy je výroba FVE minimální (pouhých $43{,}0\,\text{kWh}$ v lednu), se algoritmus soustředí na ochranu před nákupem ve špičkách a nákup v nejlevnějších nočních hodinách. Export do sítě je v zimě téměř nulový, protože baterie slouží výhradně pro pokrytí domácnosti. I přesto úspora v každém zimním měsíci činí přibližně $200\,\text{Kč}$.
	\end{itemize}

	\paragraph{3. Šetrnost k životnosti akumulátoru:}
	Za celých 17 měsíců provedl akumulátor pouze $234{,}1$ plných ekvivalentních cyklů, což odpovídá průměru přibližně \textbf{14{,}6 cyklu za měsíc}. Algoritmus díky podmínce Hurdle Rate neprovádí zbytečné mikrocyklování při malých cenových rozdílech. Při garantované životnosti $\ge 6000$ cyklů použitých článků $\text{LiFePO}_4$ představuje toto tempo opotřebení teoretickou životnost přesahující 30 let, což potvrzuje vysokou šetrnost a ekonomickou udržitelnost navrženého řešení.


	% ==========================================================================
	% LITERATURA A PRAMENY
	% ==========================================================================
	\clearpage
	\phantomsection
	\addcontentsline{toc}{chapter}{Literatura}
	\printbibliography[title={Literatura}]

	% ==========================================================================
	% PŘÍLOHY (VOLITELNÉ)
	% ==========================================================================
	\appendix
	\chapter{Struktura přiloženého zdrojového kódu}
	\label{app:struktura}

	Přiložené médium obsahuje kompletní zdrojové kódy vyvinutého systému:
	\dirtree{%
		.1 SolarEnergyOptimizer/.
		.2 src/ \DTcomment{Původní simulační moduly a historický optimalizátor}.
		.2 real\_time/ \DTcomment{Produkční služba reálného času}.
		.3 config.py \DTcomment{Kalibrace hardwaru a tarify}.
		.3 lp\_optimizer.py \DTcomment{LP optimalizační model PuLP CBC}.
		.3 controller.py \DTcomment{Řídicí logika a Battery First}.
		.3 mqtt\_service.py \DTcomment{MQTT démon s hlídacím psem (Watchdog)}.
		.3 visualizer\_realtime.py \DTcomment{Webový dashboard Streamlit}.
		.2 tests/ \DTcomment{Automatické jednotkové a integrační testy}.
		.2 bachelors\ thesis/ \DTcomment{Zdrojové texty bakalářské práce v \LaTeX u}.
	}

\end{document}
